% Hand-written appendix.  Nothing generates this file.  It is \input after
% generated/appendix.tex, which has already issued \appendix, so the \section
% below is numbered as appendix material.  It carries the scalar trace form of
% the case in Theorem~\ref{thm:quantum-coupon-collector}, which is weaker than the
% conjecture and refuted separately: the asymptotic that says why the trace
% turns negative, and the exact $n=10$ witness realizing it.
\cxappendixsection{Additional results for the quantum coupon collection conjecture}
\label{sec:qcctrace}

Theorem~\ref{thm:quantum-coupon-collector} refutes the conjecture in the Loewner order,
which is what was asked.  The failure is not confined to that ordering.  Write
\[
\Delta_n(X_1,\ldots,X_n):=\tr Q_n(X_1,\ldots,X_n)
=\sum_{\emptyset\ne S\subseteq[n]}(-1)^{|S|-1}\tr X_S^{-1},
\]
with $X_S:=\sum_{i\in S}X_i$ as in the case.  A positive definite matrix has
positive trace, so the conjecture implies $\Delta_n>0$, and
Theorem~\ref{thm:qcc-positive-five} gives $\Delta_n>0$ for every $d$ and every $n\le5$.
This scalar consequence fails too, and for a reason one can see in closed form.

Let $u_1,\ldots,u_n\in\R^3$ be such that every pair is linearly independent and every triple spans $\R^3$, and put
\[
X_i(\varepsilon):=u_iu_i^{\mathsf T}+\varepsilon I_3,\qquad\varepsilon>0,
\]
so that $X_S(\varepsilon)=\sum_{i\in S}u_iu_i^{\mathsf T}+k\varepsilon I_3$ for $|S|=k$.  A singleton leaves two regularized null directions, so $\tr X_S(\varepsilon)^{-1}=2/\varepsilon+O(1)$; a pair leaves one, so $\tr X_S(\varepsilon)^{-1}=1/(2\varepsilon)+O(1)$; and for $k\ge3$ the unregularized matrix is already positive definite, so the trace of the inverse is $O(1)$.  Only finitely many subsets occur, so the remainders sum uniformly and
\begin{equation}\label{eq:qcc-asymptotic}
\Delta_n\bigl(X_1(\varepsilon),\ldots,X_n(\varepsilon)\bigr)
=\frac1\varepsilon\Bigl(2n-\tfrac12\binom n2\Bigr)+O(1)
=\frac{n(9-n)}{4\varepsilon}+O(1).
\end{equation}
The coefficient is negative exactly for $n\ge10$, which the next theorem realizes with an exact witness.

\begin{theorem}[\statusfalse: the trace consequence]\label{thm:quantum-coupon-collector-trace}
There are $d=3$, $n=10$ and $X_1,\ldots,X_{10}\in\mathbf{S}_{++}^{3}$ with $\Delta_{10}(X_1,\ldots,X_{10})<0$.  A fortiori $Q_{10}\succ0$ fails.
\end{theorem}
\begin{proof}
Take $d=3$, $n=10$, $\varepsilon=1/10000$, and, for $t=0,1,\ldots,9$,
\[
u_t=(1,t,t^2)^{\mathsf T},\qquad X_{t+1}=u_tu_t^{\mathsf T}+\varepsilon I_3.
\]
Every $X_i$ is strictly positive definite, and every three of the $u_t$ are independent because their determinant is a nonzero Vandermonde product; the matrices are noncommuting, the $(1,2)$ entry of $[X_1,X_2]$ being $1$.  Exact rational evaluation of all $2^{10}-1=1023$ nonempty-subset terms gives
\[
-13901<\Delta_{10}(X_1,\ldots,X_{10})<-13900.
\]
The supplied verifier computes this with exact rational arithmetic and records the full numerator and denominator in \texttt{artifacts/certificate.json}.
\end{proof}

\begin{remark}[Scope]
The record for the trace form is coarser than for the conjecture itself: true
for $n\le5$ by Theorem~\ref{thm:qcc-positive-five}, false at $n=10$, and undecided for
$6\le n\le9$.  Equation \eqref{eq:qcc-asymptotic}, whose coefficient vanishes at
$n=9$, says nothing there, since that construction is only one family; and a
$Q_n$ that is not positive semidefinite, such as the one of
Theorem~\ref{thm:quantum-coupon-collector}, can still have positive trace.
\end{remark}

%%% Local Variables:
%%% mode: LaTeX
%%% TeX-master: "main"
%%% End:
